
% don't go into too much detail here.
\documentclass[../thesis.tex]{subfiles}
% this is kind  of rough here.
% Operating system
%     - describe what an operating system is

% Monolithic operating system 
%     - describe what a monolithic operating system is

% Microkernel 
%     - describe what a microkernel is, what a \MKBS would look Like, examples
    
% Hybrid kernels
%     - describe hybrid kernels
%     - subsystems, ipc, etc...

% IPC

% Access Control
%  - capabilities
%  - whatever linux uses.

% Formal verification

\begin{document}
\section{Operating System}
\begin{figure}[H]
    \centering
    \subfile{diagrams/operating_system.tex}
    \caption{Operating System in a System}
\end{figure}
The operating system is the lowest level of software directly interfacing with hardware and providing user applications a general, usable interface to operate upon. Essentially, it is an abstract machine and a resource manager working with \texttt{policies} to multiplex and allocate hardware resources through a simplified system call interface. \citep{Tanenbaum_Bos:mos5.pdf}. 

\subsection{Kernel}
A key component of any operating system is the operating system kernel. Although kernels are often referenced synonymously with operating systems, there is an important distinction to be made. Whilst parts of the operating system may run on less privileged modes, the kernel is specifically the part of the operating system which runs on the highest level of privilege. This means that the kernel has complete control and access over all computer resources in the system. Regarding the kernel, there are several designs, each having their advantages and disadvantages.
\begin{figure}[H]
    \centering
    \includegraphics{monolithic.png}
    \caption{Types of Kernels}
    \TODO{do stuff here a bout this img. use the svg somehow}
    \TODO{credit: By Golftheman - http://en.wikipedia.org/wiki/Image:OS-structure.svg, Public Domain, https://commons.wikimedia.org/w/index.php?curid=4397379}
\end{figure}

\section{Monolithic Kernel}
A monolithic kernel is a kernel design which has most if not all the operating system components inside the kernel. This usually means that all of these components run in one large address space and the kernel presents a unified system call interface to applications to access hardware resources without having to know specific details. A key advantage of a monolithic kernel is that since everything is inside the kernel, a system call only requires a mode switch into kernel mode unlike the equivalent procedure in a microkernel which requires a context switch. Linux is an example of a monolithic kernel, with all operating system components running in kernel mode. Monolithic kernels like Linux have adopted design elements from microkernels such as loadable kernel modules, which are modules that can be loaded for the operating system during runtime. It is important however to note the key difference being that loadable kernel modules run in kernel mode instead of user mode as would be the case on microkernel systems.

\section{Microkernel}
The converse of the monolithic kernel design is a microkernel which operates on Liedtke's principle of minimality \textit{`A concept is tolerated inside the $\mu$-kernel only if moving it outside the kernel, i.e., permitting competing implementations, would prevent the implementation of the system's required functionality'} \citep{Liedtke_95}. Essentially if a component is not strictly required to be inside the kernel for the system to function, then it should be implemented outside the kernel, i.e. in userland. 

A recent development is

\TODO{Probably talk about microkernel design, can quote liedtke on microkernel construction, caps and IPC etc. also talk about disadvantages and chen and bershad}
\TODO{also talk about HM as a general purpose operating system.}
\section{Hybrid Kernel}
\begin{figure}[H]
    \centering
    \includegraphics{Windows_2000_architecture.png}
    \caption{\TODO{} By The original uploader was Grm wnr at English Wikipedia.Later versions were uploaded by Xyzzy n at en.wikipedia. - Transferred from en.wikipedia to Commons., CC BY-SA 3.0, https://commons.wikimedia.org/w/index.php?curid=2242280}
    \label{Windows NT Architecture} \TODO{I need to do the svg file instead ideally.}
    
\end{figure}
Many modern operating systems are designed with a mixture of the above design philosophies. This includes Windows NT which implements certain operating system components in user space including \texttt{environment subsystems}, \texttt{server processes}, \texttt{system support processes} and \texttt{NTDLL.DLL} \TODO{cite something and etc. also do better desc of windows.}. 
and the XNU kernel which is based on the Mach microkernel as the core, including the \gls{ipc} infrastructure of Mach as a key component of the \texttt{kernel environment}, with the BSD layer being part of userland providing the application interface \citep{Apple:Kernel_Prog}. 

\section{Djawula}
Currently, monolithic and hybrid kernels dominate the general purpose operating system space and this is despite the software technologically superior design of microkernels \citep{Liedtke_95}. Even the only exception HongMeng includes many design compromises in order to achieve higher performance, to the point that one may argue that HongMeng is no longer a microkernel but instead a hybrid kernel. Having demonstrated that LionsOS can achieve higher performance than Linux on embedded workflows \citep{Heiser_VCJGNLDZAZWPDB_25}, and with a well-designed microkernel (seL4) as the kernel, is it possible to create a performant operating system?

\end{document}